\documentclass[11pt]{article}

\usepackage[T1]{fontenc}
\usepackage{lmodern}
\usepackage[margin=1in]{geometry}
\usepackage{amsmath,amssymb,amsthm,mathtools}
\usepackage[hidelinks]{hyperref}
\usepackage{microtype}

\newtheorem{theorem}{Theorem}
\newtheorem{lemma}{Lemma}
\theoremstyle{remark}
\newtheorem*{remark}{Remark}

\newcommand{\M}{\mathcal M}
\newcommand{\N}{\mathcal N}
\newcommand{\HH}{H_1^c}
\newcommand{\Tr}{\operatorname{Tr}}
\newcommand{\ind}{\mathbf 1}
\newcommand{\wstar}{\mathrm{w}^*}
\DeclareMathOperator{\dist}{dist}

\title{Weak Compactness in the Column Martingale Hardy Space}
\author{}
\date{}

\begin{document}
\maketitle

Let $(\M,\tau)$ be a finite von Neumann algebra equipped with a faithful
normal tracial state, so that $\tau(1)=1$. Let $(\M_n)_{n\geq 1}$ be an
increasing filtration with trace-preserving conditional expectations
$\mathcal E_n$. For a martingale $x=(x_n)_{n\geq 1}$, put
\[
  dx_n=x_n-x_{n-1},\qquad x_0=0,\qquad
  S_c(x)=\left(\sum_{n\geq 1}dx_n^*dx_n\right)^{1/2}.
\]
The column Hardy space $\HH=\HH(\M)$ is the completion of finite
martingales under the norm
\[
  \|x\|_{\HH}=\|S_c(x)\|_1.
\]
A family $\mathcal F\subset L_1(\M)_+$ is called \emph{uniformly
integrable} if it is norm bounded and
\[
  \lim_{R\to\infty}\sup_{h\in\mathcal F}
  \tau\!\left(h\,\ind_{(R,\infty)}(h)\right)=0.
\]
For a finite trace, the displayed condition already implies norm
boundedness. All weak topologies below are Banach-space weak topologies.

\begin{theorem}\label{thm:main}
A subset $K\subset \HH$ is relatively weakly compact if and only if
\[
  S_c(K)=\{S_c(x):x\in K\}
\]
is uniformly integrable in $L_1(\M)$.
\end{theorem}

We first record the consequence of Akemann's weak compactness criterion
that will be used in the proof. The criterion says that a bounded subset
$A\subset L_1(\N,\nu)$ is relatively weakly compact if and only if
\[
  \lim_{m\to\infty}\sup_{a\in A}\|q_m a q_m\|_1=0
\]
for every decreasing sequence of projections $q_m\downarrow 0$ in $\N$.
Here $(\N,\nu)$ may be semifinite; see \cite[Definition~2.4 and the
subsequent remarks]{Randrianantoanina}.

\begin{lemma}\label{lem:small-trace}
Let $(\N,\nu)$ be a semifinite von Neumann algebra with a faithful normal
semifinite trace. If $A\subset L_1(\N,\nu)$ is relatively weakly compact,
then
\begin{equation}\label{eq:small-trace}
  \lim_{\delta\downarrow 0}
  \sup\bigl\{\|qaq\|_1:
      a\in A,\ q\in\N\text{ a projection},\ \nu(q)\leq\delta\bigr\}=0.
\end{equation}
\end{lemma}

\begin{proof}
If \eqref{eq:small-trace} fails, there exist $\varepsilon>0$,
$a_j\in A$, and projections $r_j\in\N$ such that
\[
  \nu(r_j)\leq 2^{-j},
  \qquad \|r_j a_j r_j\|_1\geq\varepsilon
  \quad (j\geq 1).
\]
Define $q_m=\bigvee_{j\geq m}r_j$. Then $q_m$ decreases and
\[
  \nu(q_m)\leq\sum_{j\geq m}2^{-j}\longrightarrow 0.
\]
Faithfulness and normality of $\nu$ imply $q_m\downarrow 0$.
Since $r_m\leq q_m$,
\[
  \|q_m a_m q_m\|_1
  \geq\|r_m(q_m a_m q_m)r_m\|_1
  =\|r_m a_m r_m\|_1
  \geq\varepsilon.
\]
This contradicts Akemann's criterion.
\end{proof}

For completeness, we also recall a standard weak compactness
approximation principle. If $E$ is a Banach space and a bounded set
$A\subset E$ satisfies
\begin{equation}\label{eq:approximation-principle}
  A\subset W_\varepsilon+\varepsilon B_E
\end{equation}
for every $\varepsilon>0$, where $W_\varepsilon\subset E$ is weakly
compact, then $A$ is relatively weakly compact.
Indeed, identify $E$ with its canonical image in $E^{**}$ and let $C$
be the weak-star closure of $A$ in $E^{**}$. The set
$W_\varepsilon+\varepsilon B_{E^{**}}$ is weak-star compact, so
\eqref{eq:approximation-principle} implies
\[
  C\subset W_\varepsilon+\varepsilon B_{E^{**}}.
\]
Consequently $\dist(z,E)\leq\varepsilon$ for every $z\in C$ and every
$\varepsilon>0$. Since $E$ is norm closed in $E^{**}$, we have $C\subset E$.
Thus $C$ is weakly compact in $E$.

\begin{proof}[Proof of Theorem~\ref{thm:main}]
We may assume $K\neq\varnothing$. Set
\[
  \N=\M\bar\otimes B(\ell_2),\qquad
  \nu=\tau\otimes\Tr,\qquad
  p=1\otimes e_{11},\qquad
  E=L_1(\N,\nu)p,
\]
where $(e_{ij})_{i,j\geq 1}$ are the standard matrix units.
In particular, $\nu(p)=1$. The column map
\[
  J:\HH\longrightarrow E,\qquad
  Jx=\sum_{n\geq 1}dx_n\otimes e_{n1},
\]
is an isometric linear embedding. More precisely,
\begin{equation}\label{eq:column-map}
  |Jx|=S_c(x)\otimes e_{11},
  \qquad
  \|Jx\|_1=\|x\|_{\HH}.
\end{equation}
The map is first defined on finite martingales and then extended by
completion. Its image is a norm-closed, hence weakly closed, subspace
of $E$. Therefore $K$ is relatively weakly compact in $\HH$ if and only
if $A=J(K)$ is relatively weakly compact in $E$, or equivalently in
$L_1(\N,\nu)$.

\medskip
\noindent\emph{Necessity.}
Suppose that $K$ is relatively weakly compact. Then $A$ is relatively
weakly compact and
\[
  C_0=\sup_{a\in A}\|a\|_1<\infty.
\]
Fix $a=u|a|\in A$ and $R>0$. Let
\[
  e=\ind_{(R,\infty)}(|a|),\qquad
  f=ueu^*,\qquad q=e\vee f.
\]
Since $e$ is dominated by the support of $|a|$, the operator $f$ is a
projection equivalent to $e$. Hence
\begin{equation}\label{eq:tail-support}
  \nu(q)\leq\nu(e)+\nu(f)
  =2\nu(e)\leq\frac{2C_0}{R}.
\end{equation}
Moreover, $fae=ae$, and thus
\begin{equation}\label{eq:tail-compression}
  \nu\!\left(|a|\,\ind_{(R,\infty)}(|a|)\right)
  =\|fae\|_1
  =\|f(qaq)e\|_1
  \leq\|qaq\|_1.
\end{equation}
Applying Lemma~\ref{lem:small-trace} to
\eqref{eq:tail-support}--\eqref{eq:tail-compression}, we obtain
\[
  \lim_{R\to\infty}\sup_{a\in A}
  \nu\!\left(|a|\,\ind_{(R,\infty)}(|a|)\right)=0.
\]
Together with \eqref{eq:column-map}, this proves that $S_c(K)$ is
uniformly integrable.

\medskip
\noindent\emph{Sufficiency.}
Suppose that $S_c(K)$ is uniformly integrable. For $a\in A$ and $R>0$,
define
\[
  a^{(R)}=a\,\ind_{[0,R]}(|a|).
\]
The truncated operator still has right support dominated by $p$.
Since $\nu(p)=1$,
\[
  \|a^{(R)}\|_2^2
  =\nu\!\left(|a|^2\ind_{[0,R]}(|a|)\right)
  \leq R^2\nu(p)=R^2.
\]
Also,
\begin{equation}\label{eq:uniform-truncation}
  \begin{aligned}
  \eta(R)
  &:=\sup_{a\in A}\|a-a^{(R)}\|_1\\
  &=\sup_{a\in A}
     \nu\!\left(|a|\,\ind_{(R,\infty)}(|a|)\right)
    \longrightarrow 0
    \qquad (R\to\infty).
  \end{aligned}
\end{equation}
The natural inclusion
\[
  L_2(\N,\nu)p\longrightarrow E
\]
is bounded: by the noncommutative H\"older inequality, for $b=bp$,
\[
  \|b\|_1=\|bp\|_1
  \leq\|b\|_2\|p\|_2=\|b\|_2.
\]
Since $L_2(\N,\nu)p$ is a Hilbert space, the image $W_R$ of its closed
ball of radius $R$ is weakly compact in $E$. We have
\[
  A\subset W_R+\eta(R)B_E.
\]
The approximation principle above and \eqref{eq:uniform-truncation}
show that $A$ is relatively weakly compact in $E$. It follows from
the closedness of $J(\HH)$ and the isometry of $J$ that $K$ is relatively
weakly compact in $\HH$.
\end{proof}

\begin{remark}
The spectral truncations $a^{(R)}$ need not belong to $J(\HH)$.
This causes no difficulty: weak compactness is first established in the
ambient column space $E$, and then transferred to the closed subspace
$J(\HH)$.
\end{remark}

\begin{remark}[On the finite-trace hypothesis]
The theorem is stated for a finite tracial probability space. It does
not extend verbatim to arbitrary semifinite algebras.

For example, equip $B(\ell_2)$ with its usual trace and use the constant
filtration $\M_n=B(\ell_2)$. For the single-step martingales with
$x^{(k)}_n=e_{1k}$ for all $n\geq 1$, the set
$\{x^{(k)}:k\geq 1\}$ is relatively weakly compact in $H_1^c$:
the row space spanned by the $e_{1k}$ is isometric to $\ell_2$, and
$e_{1k}\to 0$ weakly in $S_1$. However,
\[
  S_c(x^{(k)})=e_{kk}.
\]
This family is not uniformly integrable in the projection-decrease
sense. Indeed, for $q_m=\sum_{k\geq m}e_{kk}\downarrow 0$,
\[
  \sup_{k\geq 1}\|q_m e_{kk}q_m\|_1=1
  \qquad (m\geq 1).
\]
Furthermore, spectral-tail decay alone does not characterize uniform
integrability in this general semifinite setting.
\end{remark}

\begin{thebibliography}{9}

\bibitem{Randrianantoanina}
N.~Randrianantoanina,
\emph{Sequences in non-commutative $L^p$-spaces},
J.~Operator Theory \textbf{48} (2002), 255--272.
Preprint: \href{https://arxiv.org/abs/math/0004144}{arXiv:math/0004144}.

\end{thebibliography}

\end{document}
